% Part: first-order-logic
% Chapter: sequent-calculus
% Section: proving-things-quant

\documentclass[../../../include/open-logic-section]{subfiles}

\begin{document}

\olfileid[pt]{fol}{seq}{prq}

\olsection{\usetoken{P}{derivation} com Quantificadores}

\begin{ex}
Dê uma $\Log{LK}$-!!{derivation} do sequente $\lexists[x][\lnot !A(x)]
\Sequent \lnot \lforall[x][!A(x)]$.

Ao lidar com quantificadores, temos de cuidar para não violar a
condição de variável própria, e às vezes isso exige que brinquemos com
a ordem de execução de certas inferências. Em geral, ajuda tentar
cuidar primeiro das regras sujeitas à condição de variável própria
(elas ficarão mais abaixo na prova finalizada). Além disso, é uma boa
ideia tentar olhar adiante e adivinhar como o sequente inicial poderia
ser. No nosso caso, ele terá de ser algo como $!A(a) \Sequent !A(a)$.
Isso significa que, ao ``reverter'' as regras de quantificador,
teremos de escolher o mesmo termo---que chamaremos de $a$---tanto para
a regra $\lforall$ quanto para a regra $\lexists$. Se escolhêssemos
termos diferentes para cada regra, acabaríamos com algo como $!A(a)
\Sequent !A(b)$, que, é claro, não é derivável.

Começando como de costume, escrevemos
\begin{prooftree}
\AxiomC{}
\UnaryInf$\lexists[x][\lnot !A(x)] \fCenter \lnot \lforall[x][!A(x)]$
\end{prooftree}
Poderíamos executar a regra \LeftR{\exists} ou a regra \RightR{\lnot}.
Como a regra \LeftR{\exists} está sujeita à condição de variável
própria, é uma boa ideia cuidar dela mais cedo do que mais
tarde, então faremos essa primeiro.
\begin{prooftree}
\AxiomC{}
\UnaryInf$ \lnot !A(a) \fCenter \lnot \lforall[x][!A(x)]$
\RightLabel{\LeftR{\lexists}}
\UnaryInf$ \lexists[x][\lnot !A(x)] \fCenter \lnot \lforall[x][!A(x)]$
\end{prooftree}
Aplicando as regras \LeftR{\lnot} e \RightR{\lnot} de trás para frente, obtemos
\begin{prooftree}
\AxiomC{}
\UnaryInf$\lforall[x][!A(x)] \fCenter !A(a)$
\RightLabel{\LeftR{\lnot}}
\UnaryInf$\lnot !A(a), \lforall[x][!A(x)] \fCenter $
\RightLabel{\LeftR{\Exchange}}
\UnaryInf$\lforall[x][!A(x)], \lnot !A(a) \fCenter $
\RightLabel{\RightR{\lnot}}
\UnaryInf$ \lnot !A(a) \fCenter \lnot \lforall[x] !A(x)$
\RightLabel{\LeftR{\lexists}}
\UnaryInf$ \lexists[x] \lnot !A(x) \fCenter \lnot \lforall[x] !A(x)$
\end{prooftree}
Neste ponto, nossa única opção é executar a regra \LeftR{\forall}.
Como essa regra não está sujeita à restrição de variável própria,
estamos livres. Lembre-se, queremos tentar obter um sequente
inicial (da forma $!A(a) \Sequent !A(a)$), então devemos escolher $a$
como nosso argumento para $!A$ ao aplicarmos a regra.
\begin{prooftree}
\Axiom$!A(a) \fCenter !A(a)$
\RightLabel{\LeftR{\lforall}}
\UnaryInf$\lforall[x][!A(x)] \fCenter !A(a)$
\RightLabel{\LeftR{\lnot}}
\UnaryInf$\lnot !A(a), \lforall[x][!A(x)] \fCenter $
\RightLabel{\LeftR{\Exchange}}
\UnaryInf$\lforall[x][!A(x)], \lnot !A(a) \fCenter $
\RightLabel{\RightR{\lnot}}
\UnaryInf$ \lnot !A(a) \fCenter \lnot \lforall[x][!A(x)]$
\RightLabel{\LeftR{\lexists}}
\UnaryInf$ \lexists[x][ \lnot !A(x)] \fCenter \lnot \lforall[x][!A(x)]$
\end{prooftree}
É importante, especialmente ao lidar com quantificadores, verificar com
cuidado neste ponto que a condição de variável própria não foi
violada. Como a única regra que aplicamos sujeita à condição de
variável própria foi \LeftR{\exists}, e a variável própria~$a$
não ocorre em seu sequente inferior (o sequente final), esta é uma
!!{derivation} correta.
\end{ex}

\begin{prob}
Dê !!{derivation}s dos seguintes sequentes:
\begin{enumerate}
\item $\Sequent (\lforall[x][!A(x)] \land \lforall[y][!B(y)]) \lif
\lforall[z][(!A(z) \land !B(z))]$.
\item $\Sequent (\lexists[x][!A(x)] \lor \lexists[y][!B(y)]) \lif
\lexists[z][(!A(z) \lor !B(z))]$.
\item $\lforall[x][(!A(x) \lif !B)] \Sequent \lexists[y][!A(y)] \lif !B$.
\item $\lforall[x][\lnot !A(x)] \Sequent \lnot\lexists[x][!A(x)]$.
\item $\Sequent \lnot\lexists[x][!A(x)] \lif \lforall[x][\lnot !A(x)]$.
\item $\Sequent \lnot\lexists[x][\lforall[y][((!A(x,y) \lif \lnot
!A(y,y)) \land (\lnot !A(y,y) \lif !A(x,y)))]]$.
\end{enumerate}
\end{prob}

\begin{prob}
Dê !!{derivation}s dos seguintes sequentes:
\begin{enumerate}
\item $\Sequent \lnot\lforall[x][!A(x)] \lif \lexists[x][\lnot!A(x)]$.
\item $(\lforall[x][!A(x)] \lif !B) \Sequent \lexists[y][(!A(y) \lif !B)]$.
\item $\Sequent \lexists[x][(!A(x) \lif \lforall[y][!A(y)])]$.
\end{enumerate}
(Todos esses requerem a regra $\RightR{\Contraction}$.)
\end{prob}

\end{document}
